\section*{Bab \ref{ch:b2:eruptive-styles} --- Erupsi Gunung Api}

\begin{solution}{exo:b2:eruptive-styles:1}
(a) $\qty{1.3}{km^3} = \qty{1.3e9}{m^3}$, antara $10^9$ dan $10^{10}$: VEI~5.
(b) $\qty{2.9e8}{m^3}$: VEI~4. (c) $\qty{2e6}{m^3}$, antara $10^6$ dan $10^7$:
VEI~2. (d) VEI~6 dimulai pada \qty{e10}{m^3}, sehingga volume total \omterm{def:b2:eruptive-styles:tephra}{tefra} sedikitnya
\qty{10}{km^3}; ``lebih dari \qty{5}{km^3}'' adalah batas bawah, bukan kontradiksi.
\end{solution}

\begin{solution}{exo:b2:eruptive-styles:2}
(a) Abu (lebih halus daripada \qty{2}{mm}). (b) \omterm{def:b2:eruptive-styles:tephra}{Lapili} (\qtyrange{2}{64}{mm}).
(c) Blok: lebih besar daripada \qty{64}{mm}, bersudut, padat ketika terlontar. (d)
\omterm{def:b2:eruptive-styles:tephra}{Bom vulkanik}, bom gelendong: lebih besar daripada \qty{64}{mm} dan terbentuk selagi lebur
saat melayang.
\end{solution}

\begin{solution}{exo:b2:eruptive-styles:3}
(a) $0.304\times(10^6)^{0.241} = 0.304\times27.9 = \qty{8.5}{km}$;
$0.304\times(10^8)^{0.241} = 0.304\times84.7 = \qty{25.8}{km}$.
(b) $Q = (35/0.304)^{1/0.241} = 115.1^{4.149} = \qty{3.6e8}{kg/s}$, dan $I =
8.55 + 3 = 11.6$. (c) $Q \propto H^{4.15}$, sehingga tinggi yang \qty{20}{\percent} terlalu tinggi
menaksir $Q$ terlalu besar dengan faktor $1.2^{4.15} = 2.1$.
\end{solution}

\begin{solution}{exo:b2:eruptive-styles:4}
$V_{\mathrm{DRE}} = 0.29\times1000/2500 = \qty{0.116}{km^3}$; $m =
\num{0.29e9}\times1000 = \qty{2.9e11}{kg}$; $M = \log_{10}(\num{2.9e11}) - 7 = 11.46
- 7 = 4.46$.
\end{solution}

\begin{solution}{exo:b2:eruptive-styles:5}
$k = \ln(40/10)/(15 - 5) = 1.386/10 = \qty{0.139}{km^{-1}}$; $T_0 = 40\,e^{0.139\times5}
= 40\times2.00 = \qty{80}{cm}$; $\ln2/k = \qty{5.0}{km}$. Dengan $T_0 =
\qty{8.0e-4}{km}$: $V = 2\times\num{8.0e-4}/0.139^2 = \qty{0.083}{km^3} =
\qty{8.3e7}{m^3}$, VEI~3.
\end{solution}

\begin{solution}{exo:b2:eruptive-styles:6}
(a) $\int_0^x T_0e^{-ky}2y\,dy = 2T_0\big[-\tfrac{y}{k}e^{-ky} - \tfrac{1}{k^2}e^{-ky}\big]_0^x
= \tfrac{2T_0}{k^2}[1 - (1 + kx)e^{-kx}]$. (b) Dengan iterasi $u = \ln(2(1 + u))$ dari
$u = 1$: $1.386$, $1.564$, $1.635$, $1.662$, $1.672$, \dots, yang konvergen ke
$u = 1.678$, dengan $(1 + u)e^{-u} = 2.678\times0.1867 = 0.500$. (c) $x_{1/2} =
1.678/0.139 = \qty{12.1}{km}$, di dalam isopak terluar ($\sqrt A =
\qty{15}{km}$); namun volume di luar isopak itu adalah $(1 + 2.09)e^{-2.09} =
\qty{38}{\percent}$ dari total. Hampir dua perlima endapan terletak di tempat yang
tidak diukur sama sekali: volumenya bertumpu pada ekstrapolasi.
\end{solution}

\begin{solution}{exo:b2:eruptive-styles:7}
(a) $p_s = (5.0/0.444)^2 = \qty{127}{MPa}$, pada $127/0.025 = \qty{5.1}{km}$.
(b) $x = (5.0 - 0.444\sqrt{10})/100 = (5.0 - 1.404)/100 = 0.0360$;
$V_g = 0.0360\times8.314\times1123/(0.018015\times10^7) = \qty{1.86e-3}{m^3}$;
$V_m = 0.964/2300 = \qty{4.19e-4}{m^3}$; $\phi = 1.86/(1.86 + 0.419) = 0.82$.
(c) Lelehan basal ribuan hingga miliaran kali kurang kental daripada lelehan
riolit (\cref{ch:b2:magma-emplacement}): gelembungnya naik,
bergabung, dan lolos, dan gas pergi melalui buih yang permeabel atau di lubang
kepundan, sehingga \omterm{def:b2:magma-emplacement:magma}{magma} melepaskan gas dalam sistem terbuka alih-alih mempertahankan gasnya;
fragmentasinya, jika terjadi, terbatas pada air mancur di dekat lubang kepundan. Kenaikan yang lebih lambat
memberikan lebih banyak waktu untuk proses ini.
\end{solution}

\begin{solution}{exo:b2:eruptive-styles:8}
(a) $u_z(0) = 0.75\times10^7/(\pi\times(5000)^2) = \qty{0.095}{m} = \qty{95}{mm}$.
(b) $z_0 = 3.8/0.766 = \qty{5.0}{km}$. (c) $(1 + 1)^{-3/2} = 0.35$ dari
maksimum.
\end{solution}

\begin{solution}{exo:b2:eruptive-styles:9}
Jangkauan terbesar dicapai pada \ang{45}: $R = v^2/g = 100^2/9.81 = \qty{1.0}{km}$, dan
$200^2/9.81 = \qty{4.1}{km}$. Hambatan udara memperpendek lintasan, sehingga nilai-nilai ini adalah
batas atas, dekat dengan kenyataan untuk blok besar dan terlalu jauh untuk blok kecil.
Karena $R = v^2/g$, zona bahaya bertambah sebagai kuadrat kecepatan lontar:
menggandakan kecepatan ledakan mengempatkalikan jari-jarinya.
\end{solution}

\begin{solution}{exo:b2:eruptive-styles:10}
(a) Dengan $T = T_0e^{-k_1x}$ untuk $x < X$ dan $T = T_Xe^{-k_2(x - X)}$ di luarnya, $T_X =
T_0e^{-k_1X}$ (kontinuitas), gunakan $\int_0^X ye^{-ky}dy = \frac{1}{k^2} - e^{-kX}\big(\frac{X}{k} +
\frac{1}{k^2}\big)$ dan $\int_X^\infty ye^{-ky}dy = e^{-kX}\big(\frac{X}{k} +
\frac{1}{k^2}\big)$:
bagian proksimal memberikan $2T_0\big[\frac{1}{k_1^2} - e^{-k_1X}\frac{k_1X + 1}{k_1^2}\big]$ dan
bagian distal $2T_Xe^{k_2X}\int_X^\infty ye^{-k_2y}dy = 2T_X\frac{k_2X + 1}{k_2^2}$; dengan $T_X =
T_0e^{-k_1X}$ jumlah keduanya adalah rumus pada kotak metode.
(b) $T_0 = \qty{0.01285}{km}$: $2T_0/k_1^2 = 0.0257/0.1332 = \qty{0.193}{km^3}$;
$e^{-0.365\times8.56} = 0.0440$; $(0.138\times8.56 + 1)/0.138^2 = 114.5$ dan
$(0.365\times8.56 + 1)/0.365^2 = 30.96$; koreksinya $0.0257\times0.0440\times83.6
= \qty{0.094}{km^3}$ dan $V = \qty{0.287}{km^3}$, yaitu \qty{0.29}{km^3} yang diterbitkan.
Di luar $X$ terletak $2T_X(X/k_2 + 1/k_2^2) = 2\times\num{5.65e-4}\times114.5 =
\qty{0.129}{km^3}$, yaitu \qty{45}{\percent}.
\end{solution}

\begin{solution}{exo:b2:eruptive-styles:11}
(a) $100^{1/4} = 3.16$; hukum-hukum empiris memberikan $100^{0.241} = 3.03$ dan
$100^{0.226} = 2.83$. (b) $H \propto N^{-3/4}$: atmosfer yang terstratifikasi lebih kuat
menghentikan kolom lebih rendah; untuk $N$ yang \qty{20}{\percent} lebih besar,
$1.2^{-0.75} = 0.87$, tinggi yang \qty{13}{\percent} lebih rendah. (c) Plume yang lemah
dibengkokkan oleh angin, yang merendahkannya, sedangkan udara lembap melepaskan kalor laten di
dalamnya, yang meninggikannya; pengaruh-pengaruh ini berbobot berbeda pada erupsi kecil dan besar
dan tidak termuat dalam hukum sederhana itu. Selain itu, lajunya hanya diketahui dalam batas
faktor beberapa kali, galat pada variabel sumbu mendatar yang membiaskan
kemiringan hasil pencocokan ke arah nol.
\end{solution}

\begin{solution}{exo:b2:eruptive-styles:12}
(a) $132/400\times100 = 33$ per abad sebelum 1900; $95/125\times100 = 76$ per
abad sesudahnya. (b) $33/76 = \qty{43}{\percent}$. (c) $21/56 = \qty{38}{\percent}$:
pencatatan kurang yang sama, dengan faktor dua hingga tiga. (d) Erupsi yang lebih besar
terlihat dari tempat yang lebih jauh dan meninggalkan endapan yang lebih tebal dan luas yang bertahan terhadap erosi
lebih lama; erupsi VEI~4 di pulau terpencil pada tahun 1700 dapat lolos dari catatan, erupsi
VEI~6 hampir tidak mungkin.
\end{solution}

\begin{solution}{pb:b2:eruptive-styles:1}
\textbf{1.} $\ln T = 5.298$, $4.605$, $3.871$, $2.639$, $1.386$. Kemiringan antara
isopak-isopak yang berurutan adalah $-0.0693$, $-0.0367$, $-0.0308$, dan $-0.0313$ per km: satu
segmen terjal antara 10 dan \qty{20}{km}, satu segmen landai dari 40 hingga
\qty{120}{km}; patahannya terletak antara 20 dan \qty{40}{km}.

\textbf{2.} $k_1 = \ln(200/100)/10 = \qty{0.0693}{km^{-1}}$; $T_0 = 200\,e^{0.693} =
\qty{400}{cm}$.

\textbf{3.} $167e^{-1.244} = \qty{48.1}{cm}$, $167e^{-2.488} = \qty{13.9}{cm}$,
$167e^{-3.732} = \qty{4.00}{cm}$: dalam selisih \qty{1}{\percent} dari 48, 14, dan 4.0.

\textbf{4.} $400e^{-0.0693X} = 167e^{-0.0311X}$ memberi $X = \ln(400/167)/(0.0693 -
0.0311) = 0.873/0.0382 = \qty{22.9}{km}$.

\textbf{5.} $\ln2/0.0693 = \qty{10.0}{km}$ dan $\ln2/0.0311 = \qty{22.3}{km}$.

\textbf{6.} $T_0 = \qty{0.0040}{km}$: $V = 2\times0.0040/0.0693^2 =
\qty{1.67}{km^3}$.

\textbf{7.} $e^{-0.0693\times22.9} = 0.205$; $(0.0311\times22.9 + 1)/0.0311^2 = 1770$;
$(0.0693\times22.9 + 1)/0.0693^2 = 539$; koreksi $2\times0.0040\times0.205\times1231 =
\qty{2.02}{km^3}$; $V = 1.67 + 2.02 = \qty{3.68}{km^3}$, dan koreksinya
\qty{55}{\percent} dari volume itu.

\textbf{8.} $3.68/1.67 = 2.2$. Abu distal yang halus itu tipis, mudah tererosi,
terolah ulang, atau jatuh ke laut; pada endapan tua abu itu sering hilang atau tidak pernah
dipetakan, dan volumenya kemudian ditaksir terlalu rendah dengan faktor dua atau lebih.

\textbf{9.} $\qty{3.68e9}{m^3}$ terletak antara $10^9$ dan $10^{10}$: VEI~5.

\textbf{10.} $V_{\mathrm{DRE}} = 3.68\times900/2400 = \qty{1.38}{km^3}$.

\textbf{11.} $m = \num{3.68e9}\times900 = \qty{3.31e12}{kg}$.

\textbf{12.} $M = \log_{10}(\num{3.31e12}) - 7 = 12.52 - 7 = 5.52$.

\textbf{13.} $3.68/1.3 = 2.8$ kali.

\textbf{14.} Kolom yang lebih tinggi daripada \qty{25}{km}, dan gaya Plini.

\textbf{15.} $Q = \num{3.31e12}/36000 = \qty{9.2e7}{kg/s}$; $I = 7.96 + 3 = 11.0$.

\textbf{16.} $0.304\times(\num{9.2e7})^{0.241} = 0.304\times83.0 = \qty{25.2}{km}$.

\textbf{17.} $0.345\times(\num{9.2e7})^{0.226} = 0.345\times63.1 = \qty{21.8}{km}$.

\textbf{18.} Nilai terukur \qty{31}{km} \qty{23}{\percent} di atas ramalan
pertama dan \qty{42}{\percent} di atas ramalan kedua, masih di dalam sebaran
\qtyrange{53}{60}{\percent} hukum-hukum itu: konsisten.

\textbf{19.} $Q = (31/0.304)^{1/0.241} = 102.0^{4.149} = \qty{2.2e8}{kg/s}$, 2.3 kali
laju rata-rata. Klas merekam kolom tertinggi, pada puncak
erupsi, sedangkan pertanyaan~15 merata-ratakan atas seluruh fase; dan sebagian massa
mungkin tidak terekam dalam endapan (abu distal, \omterm{def:b2:eruptive-styles:pdc}{aliran piroklastik}).

\textbf{20.} Ya: 22 hingga \qty{31}{km}, jauh di atas tropopause pada \qty{16}{km}.
Gas yang terbawa ke stratosfer menyebar ke seluruh dunia; gas kuncinya adalah
belerang dioksida, dan dua puluh juta ton yang dimasukkan Pinatubo ke sana pada tahun 1991
mendinginkan dunia sekitar setengah derajat selama dua tahun.

\textbf{21.} $\lambda_3 = 399/75 = 5.32$, $\lambda_4 = 95/125 = 0.760$, $\lambda_5 = 22/225 =
0.0978$, $\lambda_6 = 8/525 = 0.0152$ per tahun; logaritmanya $0.726$, $-0.119$, $-1.010$,
$-1.817$.

\textbf{22.} Garis itu memberikan $10^{0.73} = 5.37$, $10^{-0.12} = 0.759$, $10^{-0.97} =
0.107$, dan $10^{-1.82} = 0.0151$: dalam selisih \qty{1}{\percent}, \qty{0}{\percent},
\qty{10}{\percent}, dan \qty{1}{\percent} dari hitungannya.

\textbf{23.} Rekaman erupsi besar tetap lengkap lebih jauh ke belakang,
karena erupsi itu teramati dari tempat yang lebih jauh dan endapannya bertahan lebih lama; jendela
yang dimulai pada 1500 akan sangat kurang menghitung erupsi VEI~3, sedangkan jendela
yang dimulai pada 1950 hampir tidak akan memuat erupsi VEI~6.

\textbf{24.} $\lambda_5 = \num{0.107}$ per tahun, periode ulang $1/0.107 =
\num{9.3}\text{ tahun}$; $1 - e^{-1.07} = 1 - 0.343 = \qty{66}{\percent}$ dalam sepuluh tahun.

\textbf{25.} $\lambda_7 = 10^{3.28 - 5.95} = 10^{-2.67} = 0.00214$ per tahun, satu setiap
470 tahun; rekaman memberikan $2/1025 = 0.00195$ per tahun: sesuai dalam selisih
\qty{10}{\percent}, untuk titik yang tidak dipakai dalam pencocokan.

\textbf{26.} $\lambda_6 = 10^{3.28 - 5.10} = 10^{-1.82} = 0.0151$ per tahun; $\lambda t = 1.51$
dan $1 - e^{-1.51} = 1 - 0.220 = 0.780$: peluang sedikitnya satu erupsi
VEI~6 atau lebih di suatu tempat di Bumi dalam seratus tahun mendatang adalah
\textbf{\qty[text-series-to-math=true]{78}{\percent}}. Angka ini bertumpu pada laju yang tetap,
pada erupsi-erupsi yang saling independen (proses Poisson), dan pada
katalog yang lengkap di dalam jendela-jendelanya.
\end{solution}
